% method_allocation.tex
% Methods: the score gamma and the selection rules, then allocation by constrained utility maximization.
% Agent A draft, 2026-10-30. Verified against the code, which is the truth where the design notes differ:
% code/src/Allocators.jl, AllocatorsLongShort.jl, AllocatorsMeanVariance.jl (utilities, limits, solvers, benchmarks),
% code/src/HedgePool.jl (pool, down-day beta, beta cap), code/src/StudyBacktest.jl (_study_problem,
% _study_problem_columns, _study_decision, _study_lots, run_study), code/src/MarketInputs.jl (xi, g-tilde, gamma),
% code/src/SIM.jl and OnlineSIM.jl (EWLS), code/src/Covariance.jl (Sigma_g), code/src/Screens.jl (selection rules,
% bandit), config/study.toml (every setting), results/calibration_summary.toml (calibrated values).
% Notation follows methods_data.tex and method_mpc.tex: bar c is the half spread as a fraction (5 bps), Sigma_g is
% Eq. eq:mpc-cov, g_i(k) and g_M(k) are daily growth rates (Eq. eq:growth).
% The source of each number is in a comment on the line above it.

\subsection{The score and the selection rules}\label{sec:method-score}

The score combines a single-index fit of each name to the market with the market state. The daily growth rates
$g_i(k)$ of name $i$ and $g_M(k)$ of SPY (Eq.~\ref{eq:growth}, yr$^{-1}$) follow the single-index model
$g_i(k) = \alpha_i + \beta_i\, g_M(k) + \varepsilon_i(k)$~\cite{sharpe1963}, with intercept $\alpha_i$ (yr$^{-1}$),
market loading $\beta_i$ (dimensionless) and residual $\varepsilon_i(k)$ (yr$^{-1}$). We fit the model by ordinary
least squares on 2014--2024 and, during the trading window, updated $(\alpha_i, \beta_i)$ every trading day by
exponentially weighted least squares (EWLS),
\begin{equation}
(\alpha_i, \beta_i)_t = \arg\min_{a,\,b} \Big[ \sum_{k = t_0 + 1}^{t} \delta^{\,t-k}
  \big( g_i(k) - a - b\, g_M(k) \big)^2
  + \delta^{\,t - t_0} N_0\, \mathrm{E}_0 \big[ ( g_i - a - b\, g_M )^2 \big] \Big],
\label{eq:ewls}
\end{equation}
% config/study.toml [sim] ewls_half_life = 63.0, ewls_prior_weight = 63.0; SIM.jl ewls_init, ewls_update!
where $t_0$ is the last trading day before the window and $\delta = 2^{-1/63}$ is the daily decay, a half-life of
$63$ trading days. The prior term counts as $N_0 = 63$ observations, with $\mathrm{E}_0$ the expectation under the
% results/calibration_summary.toml m_m = 0.1061 (1/yr), v_m = 7.511 (1/yr^2)
2014--2024 fit and SPY's 2014--2024 mean ($0.106$~yr$^{-1}$) and variance ($7.51$~yr$^{-2}$) of daily growth; its
minimizer is the 2014--2024 fit, and its weight decays with the data's. The estimates used on day $t$ include the
growth into day $t$ and nothing later. Two market inputs come from SPY's closes. The market growth average
% config/study.toml [market] gtilde_window = 252; MarketInputs.jl gtilde_series
$\tilde g_M$ (yr$^{-1}$) is the exponential moving average of SPY's daily growth rate with a span of $252$ trading
days (smoothing factor $2/253$). The market state is
\begin{equation}
\xi_t = -G \left( \frac{\mathrm{EMA}_{21}(t)}{\mathrm{EMA}_{63}(t)} - 1 \right),
\label{eq:xi}
\end{equation}
% config/study.toml [market] xi_short = 21, xi_long = 63, xi_quantile = 0.99; MarketInputs.jl xi_gain, xi_series
where $\mathrm{EMA}_{21}$ and $\mathrm{EMA}_{63}$ are exponential moving averages of SPY's close with spans of $21$
% results/calibration_summary.toml G = 21.2315
and $63$ trading days, and the gain $G = 21.23$ is the reciprocal of the $99$th percentile of
$\lvert \mathrm{EMA}_{21}/\mathrm{EMA}_{63} - 1 \rvert$ over 2014--2024, so $\lvert\xi_t\rvert \le 1$ on $99\%$ of
those days. The state $\xi_t$ is positive exactly when the 21-day average lies below the 63-day average, as in a
% talk/data/market_state.csv: xi > 0 on 75 of 420 trading days
falling market, and it was positive on $75$ of the $420$ trading days of the window (Fig.~\ref{fig:market-state}).
The score of name $i$ is
\begin{equation}
% config/study.toml [market] beta_floor = 0.1; MarketInputs.jl gamma_weights
\gamma_i = \tanh\!\left( \frac{\alpha_i + \beta_i\, \tilde g_M}{\max(\lvert\beta_i\rvert,\, 0.1)^{\xi_t}} \right),
\label{eq:gamma}
\end{equation}
with growth rates in yr$^{-1}$ entering the hyperbolic tangent as plain numbers, so $\gamma_i \in (-1, 1)$. The
numerator is the name's expected growth rate $\mu_{g,i}$ from the single-index fit. It keeps the sign of $\beta_i$,
so a name with a negative loading is charged rather than credited for market growth when $\tilde g_M > 0$. The
denominator, the beta penalty, is positive, so $\gamma_i$ always has the sign of the expected growth rate and the
market state changes only its size. When $\xi_t > 0$ the penalty divides the expected growth of a name with
$\lvert\beta_i\rvert > 1$ by a number above one and that of a name with $\lvert\beta_i\rvert < 1$ by a number below
one, so the score favors names that move less with the market; when $\xi_t < 0$ the effect reverses and the score
favors names that move more; at $\xi_t = 0$ the score is the hyperbolic tangent of expected growth. The floor of
$0.1$ keeps the penalty finite for names with a loading near zero, and at the window's largest $\xi_t$ of $0.90$ it
% 0.1^(-0.90) = 7.94
limits their amplification to a factor of $0.1^{-0.90} \approx 7.9$. The same expected growth rate enters the
tangency rule and the controller's forecast.

% results/calibration_summary.toml universe_size = 373, 11 sectors
A selection rule chose the names of each decision from the $373$ names of the universe in $11$ sectors, using data
dated on or before the day of the pick only, and the study compared ten rules. Three picked $3$ names in each sector,
$33$ in all: \emph{$\gamma$ per sector} kept the $3$ highest $\gamma_i$ in each sector, \emph{bandit per sector} kept
the $3$ names a per-sector bandit ranked highest, and \emph{random per sector} drew $3$ names per sector uniformly at
% config/study.toml [screens] basket_K = 13
random. Five picked $K$ names from the whole universe with no sector quota, with $K = 13$ in the main comparison and
other sizes in the size comparison (Fig.~\ref{fig:S-size}). \emph{Global $\gamma$} kept the $K$ highest $\gamma_i$.
\emph{Correlation-aware $\gamma$} walked down the same ranking and accepted a name only if the Pearson correlation of
% config/study.toml [screens] corr_max = 0.7, corr_window = 252
its daily growth rates with those of every name already accepted, over the trailing $252$ trading days, was at most
$0.7$, and filled any shortfall from the ranking. \emph{Liquid then $\gamma$} kept the $K$ highest $\gamma_i$ among
% config/study.toml [screens] liquidity_shortlist = 100, liquidity_window = 63
the $100$ names with the largest median daily dollar volume (close times volume) over the trailing $63$ trading days,
\emph{random liquid 100} drew $K$ names uniformly from that shortlist, and \emph{random whole universe} drew $K$ names
uniformly from the universe. The random rules drew afresh at every pick, from a stream seeded by the run seed and the
date. The last two rules were fixed lists: the hand-picked $13$ (AAPL, MSFT, INTC, MU, AMD, NVDA, GS, BAC, WFC, C, F,
% config/study.toml [grid.fixed_lists] l8b_13, l8b_15
GM, JNJ), which the author chose in advance from sectors the author follows and without 2025 data, and the same list
with WMT and PG added. Each is one list chosen by one person, and a pick never changed it, although the hedge pool
(the stocks with zero or negative beta that join each decision when the market hedge is on) could change.

% config/study.toml [screens] bandit_*; Screens.jl MyBanditScreen, pick, train_ticker_bandit, rank_arms
The bandit rule ran one $\epsilon$-greedy bandit~\cite{lattimoreSzepesvari2020} per sector, retrained from scratch at
% results/calibration_summary.toml sector_counts: 15 (Communication Services) to 60 (Industrials)
each pick, with each of the sector's $A$ names (between $15$ and $60$) as one arm. A pull of arm $j$ drew a window
start $d$ uniformly from the $252$ trading days whose 21-day window ends within the trailing $252$ trading days, so
every window was complete by the day of the pick, and earned the excess growth rate
$x_j(d) = [\ln(S_j(d+21)/S_j(d)) - \ln(S_M(d+21)/S_M(d))]/(21\,\Delta t)$ (yr$^{-1}$) of name $j$ over SPY, with
$S_M$ SPY's close. The bandit made $\min(\max(50A, 500), 5000)$ pulls. The first pull chose an arm uniformly at
random; at pull $k > 1$ the bandit chose a uniformly random arm with probability
$\epsilon_k = \max\{0.05,\, \min[1,\, k^{-1/3} (A \ln k)^{1/3}]\}$, and otherwise the pulled arm with the highest
running-mean reward, ties going to a random rank drawn at the start. With these settings $\epsilon_k$ stayed above
% computed from the epsilon_k formula with A = 15..60 and the pull counts above: min over pulls and sectors = 0.510
$0.51$ at every pull in every sector, so the floor of $0.05$ never bound and each pull explored with probability
above one half. The rule kept the $3$ arms with the highest mean reward, unpulled arms last
(Algorithm~\ref{alg:S-bandit}). Because every arm can be scored on every window, the bandit approximates a ranking of
each sector's names by their mean 21-day growth rate in excess of SPY over the trailing year, a sector-neutral
momentum screen~\cite{jegadeesh1993}.


\subsection{Allocation by constrained utility maximization}\label{sec:method-allocation}

At each decision, made after the close of trading day $t$, the allocator chose a continuous number of shares
$n_i \ge 0$ for every name $i$ in the decision's name set $\mathcal{N}$, which held the names picked by the selection
rule and, with the market hedge on, a pool of hedge names. With $S_i$ the closing price of name $i$ (USD/share), $W$
the wealth at the close (USD) and $w_i = n_i S_i / W$ the weight of name $i$ (dimensionless), the allocator solved
\begin{equation}
\begin{aligned}
\max_{n \ge 0}\quad & U(n;\gamma) \\
\text{subject to}\quad & \textstyle\sum_{i \in \mathcal{N}} n_i S_i \le W, \qquad
  w_i \le \bar w_i \;\; (i \in \mathcal{N}), \qquad \sqrt{w^{\top} \Sigma_g\, w} \le \sigma_{\mathrm{SPY}}, \\
& \bar c\, \bar S \textstyle\sum_{i \in \mathcal{N}} \lvert n_i - n_i^{\mathrm{prev}} \rvert \le B, \qquad
  \textstyle\sum_{i \in \mathcal{N}} \beta_i w_i \le \bar\beta \;\; \text{(hedge on)},
\end{aligned}
\label{eq:allocation}
\end{equation}
where $U(n;\gamma)$ is the weighting rule's utility and $\gamma_i$ the score of name $i$ (Eq.~\ref{eq:gamma}); wealth
not invested is held as cash, which has a market loading of zero.
% config/study.toml [limits] w_max = 0.10, w_max_scale = 1.5; StudyBacktest.jl _study_problem (cap_pick, is_short)
The single-name cap was $\bar w_i = \max(0.10,\, 1.5/K)$ for a name picked by the selection rule, with $K$ the number
of picks, and $0.10$ for hedge-pool names and $130/30$ short names, so a pick could hold up to $10\%$ of wealth in the
% K = 33, 13, 5: max(0.10, 1.5/K) = 0.10, 0.1154, 0.30
$33$-name baskets, $11.5\%$ in the $13$-name baskets and $30\%$ at $K = 5$. In the volatility cap, $\Sigma_g$
(yr$^{-2}$) is the covariance of the names' daily growth rates, the single-index market term plus the sample
% config/study.toml [covariance] omega_window = 126; results/calibration_summary.toml lambda = 0.0
residual covariance over the trailing $126$ trading days (Eq.~\ref{eq:mpc-cov} with its calibrated shrinkage of
% results/calibration_summary.toml sigma_m = 2.7411 (1/yr), vol_cap_annual = 0.1727
zero), and $\sigma_{\mathrm{SPY}} = 2.741$~yr$^{-1}$ is the sample standard deviation of SPY's daily growth rate over
2014--2024, an annualized volatility of $17.3\%$; the solver imposed the cap as a second-order cone on $L^{\top} w$
% Allocators.jl add_shared_constraints!: cholesky(Sigma_g + 1e-8 I)
with $L L^{\top} = \Sigma_g + 10^{-8} I$. The turnover budget caps the half-spread cost of the share changes from
the holdings $n^{\mathrm{prev}}$ before the decision, priced at the half spread $\bar c$ of the cost model
(Eq.~\ref{eq:cost}) times the mean closing price $\bar S$ (USD/share) of the names in $\mathcal{N}$, with
% config/study.toml [costs] half_spread_bps = 5.0; [limits] turnover_frac = 0.10
$\bar c = 5 \times 10^{-4}$ and $B = 0.10\,W$ (USD). For a long-only book this budget can bind only when $\bar S$
% StudySettings.jl docstring: 5e-4 * mean(S) * 2W/min(S) <= 0.10 W unless mean(S)/min(S) > 100
exceeds $100$ times the cheapest name's price. The sum in the budget ran over the names whose shares the solver
chose; the sale of a name with $\gamma_i \le 0$ down to its floor and the exit of a held name outside $\mathcal{N}$
were not counted.
% config/study.toml [limits] s_max = 1.0 (no sector cap)
The study set no sector limit. With the hedge on, the portfolio's market loading $\sum_i \beta_i w_i$, computed with
each name's signed online loading (Eq.~\ref{eq:ewls}), was capped at
% config/study.toml [hedge] beta_ref = 1.0, kappa_beta = 1.0, beta_min = 0.0; HedgePool.jl beta_cap
$\bar\beta = \max(0,\, 1 - \xi_t)$, where $\xi_t$ is the market state (Eq.~\ref{eq:xi}). The cap equals SPY's loading
of one at $\xi_t = 0$, tightens as $\xi_t$ rises in a falling market, reaches zero at $\xi_t \ge 1$ and rises above
% talk/data/market_state.csv: min xi = -0.7999, max xi = 0.8999 over the 420 trading days
one when $\xi_t < 0$; over the trading window $\xi_t$ ranged from $-0.80$ to $0.90$, so $\bar\beta$ ranged from
$1.80$ down to $0.10$ (Fig.~\ref{fig:market-state}). The cap applied to the targets of each re-plan. Between
re-plans the shares were held, so the book's loading moved with prices and the online loadings and was not held
under the cap.

% config/study.toml [hedge] pool_H = 5, downday_lookback = 252; HedgePool.jl down_day_beta, hedge_pool
The hedge pool was refreshed at every pick of the basket. A universe name was eligible when its online market loading
was at most zero and its down-day beta was at most zero, where the down-day beta is the ordinary least squares slope,
with intercept, of the name's daily growth rate on SPY's over the days of the trailing $252$ trading days on which
SPY's growth rate was negative. The pool held up to $H = 5$ eligible names, most negative down-day beta first, and
joined every decision's name set on top of the picks, a name both picked and pooled counting once as a hedge name.
In the rules that use $\gamma$, a pool name with $\gamma_i > 0$ entered the utility like any other name, while a pool
name with $\gamma_i \le 0$ had no utility term and was a free variable above its floor, so the allocator bought it only
when the $\beta$ cap or the volatility cap made it worth its share of the budget. Because the objective is flat along
% Allocators.jl _HEDGE_TIEBREAK = 1e-6, _min_hedge, _solve_utility
such a name's weight whenever no limit binds, it subtracted a tie-break cost of $10^{-6}$ per unit of hedge weight,
and a second, lexicographic solve then fixed every other name's shares at their first-stage values and minimized the
total weight of the free hedge names under the same limits. The utility depends only on the fixed names, so the
second solve left it unchanged and held a hedge name above its floor only as far as a binding limit required; if the
second solve failed, the first-stage shares stood. With the hedge off, a decision had neither the $\beta$ cap nor the
% docs/superpowers/specs/2026-10-20-screened-utility-panel-study-design.md section 14 (known limitations)
pool and was otherwise identical. We added the $\beta$ cap and the pool to the design after inspecting how the online
loadings behaved in the February--April 2025 sell-off; every decision used trailing data only, and every result is
reported with the hedge on and off.

We compared four utilities built on the score. With $\mathcal{P} = \{ i \in \mathcal{N} : \gamma_i > 0 \}$ the names
with a positive score, they are
\begin{equation}
\begin{aligned}
U_{\mathrm{CD}} &= \sum_{i \in \mathcal{P}} \gamma_i \ln n_i, &
U_{\mathrm{CES}} &= \Big( \sum_{i \in \mathcal{P}} \gamma_i\, n_i^{\rho} \Big)^{1/\rho}, \\
U_{\mathrm{SG}} &= \sum_{i \in \mathcal{P}} \gamma_i \ln \left( n_i - b_i \right), &
U_{\mathrm{CESw}} &= \Big( \sum_{i \in \mathcal{P}} \gamma_i\, w_i^{\rho} \Big)^{1/\rho},
\end{aligned}
\label{eq:utilities}
\end{equation}
the Cobb--Douglas (CD)~\cite{cobbDouglas1928}, Stone--Geary (SG)~\cite{geary1950, stone1954} and constant elasticity
of substitution (CES)~\cite{arrowEtAl1961} utilities, the last on share counts (CES) and on weights (CESw). Both CES
% config/study.toml [allocators] ces_eta = 0.5, sg_baseline = 0.01
rules used the elasticity $\eta = 0.5$, so $\rho = 1 - 1/\eta = -1$, and SG used the baseline $b_i = 0.01\,W/S_i$
shares, a core position worth $1\%$ of wealth that SG holds in every name of $\mathcal{N}$ ($n_i \ge b_i$). Under the
budget alone each utility has a closed-form optimum (Table~\ref{tab:utilities}; derivations in
Appendix~\ref{app:allocation}). CD splits the invested wealth in proportion to $\gamma_i$. CES sets
$n_i \propto (\gamma_i/S_i)^{\eta}$, so its weights $w_i \propto \gamma_i^{\eta} S_i^{1-\eta}$ depend on each name's
share price when $\eta \ne 1$, whereas CESw sets $w_i \propto \gamma_i^{\eta}$ and does not; both equal CD at
$\eta = 1$, approach equal share counts (CES) or equal weights (CESw)~\cite{demiguelGarlappiUppal2009} as
$\eta \to 0$, and concentrate in the name with the largest $\gamma_i/S_i$ (CES) or $\gamma_i$ (CESw) as
$\eta \to \infty$. SG first buys every baseline and then splits the remaining wealth in proportion to $\gamma_i$.
With the limits of Eq.~\ref{eq:allocation} there is no closed form, and each utility becomes a convex cone
program~\cite{boydVandenberghe2004}: CD and SG use exponential cones, the two CES rules use power cones, the
volatility cap is a second-order cone and the other limits are linear. We built each program in JuMP
% code/Manifest.toml: JuMP 1.32.1, Clarabel 0.11.1, SCS 2.7.0; Allocators.jl _CLARABEL_ATTRIBUTES
$1.32.1$~\cite{lubin2023jump} and solved it with the interior-point solver Clarabel $0.11.1$~\cite{goulart2024clarabel}
(absolute and relative duality-gap tolerances $10^{-12}$, feasibility tolerance $10^{-9}$, at most $500$ iterations),
accepting its optimal and almost-optimal statuses; when Clarabel reached neither, the same program went to SCS
$2.7.0$~\cite{odonoghue2016scs} with its default settings. To check the solver against the closed forms, we allocated
% results/limits_day.csv; scripts/11_run_study.jl limits_day_step; StudyRunner.jl limits_runs; config/study.toml W0
a new book of $100{,}000$ USD on 2 January 2025 across the hand-picked $13$-name list, with the hedge off and every
limit except the budget removed (the turnover budget stayed and did not bind). Over $17$ solves (CD, and CES and CESw
% results/limits_day.csv: max |weight - closed_form| = 1.03e-5 (CES, eta = 0.1); gate tol 1e-3 (Limits.jl)
at eight values of $\eta$ from $0.05$ to $20$), the largest gap between a solver weight and its closed form was
$1.0 \times 10^{-5}$, below the tolerance of $10^{-3}$ we required.

% Allocators.jl allocate (CD, CES, CESw: pins at constraints.eps); Backtest.jl _BT_eps = 1e-3, _BT_HELD_W = 1e-3
The rules treat names with $\gamma_i \le 0$ differently. CD, CES and CESw leave them out of the utility and hold them
at a floor of $\varepsilon = 10^{-3}$ shares, which disappears when targets worth at most $0.1\%$ of wealth are set
to zero before trading, so the three rules held only names with $\gamma_i > 0$ and any hedge names the limits
required; when no name in $\mathcal{N}$ had $\gamma_i > 0$, they skipped the solve and moved the book to cash. SG
holds a name with $\gamma_i \le 0$ at its baseline $b_i$; with no positive score it held the baselines, adding the
least hedge weight that met the $\beta$ cap, and it reported the problem infeasible when the baselines exceeded the
budget or broke another limit. The baselines use $1\%$ of wealth per name, so SG cannot hold more than $100$ names.
% results/large_k/large_k_summary.csv: SG n_failed 600 of 600 decisions at K = 200 (mean invested 0), 2,435 of 2,768 at K = 100
In the large-book runs every SG decision at $K = 200$ was infeasible and the book stayed in cash, and at $K = 100$ most
SG decisions were infeasible and kept the previous book.

Three benchmarks ran under the same limits. Minimum variance (MinVar) and tangency
(Tangent)~\cite{markowitz1952, tobin1958} do not use $\gamma$ and could hold any name of $\mathcal{N}$, with
$n_i \ge 0$ and no floors. Each followed a two-fund form: it found the largest invested fraction
$s = \min(1, \max \sum_i w_i)$ that the budget, single-name, $\beta$ and turnover limits allowed without the
volatility cap, chose the mix at $\sum_i w_i = s$ that minimized $w^{\top} \Sigma_g w$ (MinVar) or maximized
$(\mu_g - g_f \mathbf{1})^{\top} w / \sqrt{w^{\top} \Sigma_g w}$ (Tangent), and scaled the mix by
$\min(1, \sigma_{\mathrm{SPY}} / \sigma_{\mathrm{mix}})$ onto the volatility cap, holding the rest in cash. Here
$\mu_{g,i} = \alpha_i + \beta_i \tilde g_M$ (yr$^{-1}$) is the expected growth rate of name $i$ from the score's inputs
and $g_f$ (yr$^{-1}$) the risk-free growth rate (Eq.~\ref{eq:riskfree}). Tangent solved its ratio as a quadratic
program after a Charnes--Cooper change of variables~\cite{charnes1962} and fell back to MinVar when no feasible mix
had positive excess growth. The $130/30$ long-short Cobb--Douglas rule (LS)~\cite{jacobsLevyStarer1998} maximized
$\sum_{\gamma_i > 0} \gamma_i \ln n_i^{+} + \sum_{\gamma_i < 0} \lvert\gamma_i\rvert \ln n_i^{-}$ over long shares
$n_i^{+}$ and short shares $n_i^{-}$ ($n_i = n_i^{+} - n_i^{-}$), so it shorted names with $\gamma_i < 0$ weighted by
% config/study.toml [allocators] ls_long = 1.30, ls_short = 0.30, ls_net = 1.00
$\lvert\gamma_i\rvert$. Long weights summed to at most $1.30$, short weights to at most $0.30$ in absolute value and
the net weight to at most $1$, so cash stayed nonnegative; the single-name cap applied to $\lvert w_i \rvert$, the
volatility and $\beta$ caps to the signed weights and the turnover budget to the signed share changes. Its long side
held the picks with $\gamma_i > 0$ and the hedge pool, and its short side held the short candidates, up to $3$ names
% config/study.toml [screens] N_per_sector = 3; Screens.jl short_candidates
per sector with the most negative $\gamma_i < 0$ outside the pool, chosen at each pick of the basket and kept only
while $\gamma_i < 0$. A hedge name was never shorted, since shorting a name with a negative loading raises the
portfolio's loading. With no short candidate LS ran as long-only CD, and with no positive score it moved to cash.

Algorithm~\ref{alg:allocation} lists one allocation decision. A decision whose program a solver certified
infeasible, that neither solver solved, or that produced non-finite shares kept every holding and placed no orders.
% results/tier1_summary.csv: n_failed 102 of n_decisions 23,050 (0.44%); results/fixed_summary.csv: 105 of 11,497
% (0.91%, the largest); results/universe_summary.csv: 171 of 121,859 (0.14%); results/gamma_seeds/gamma_seeds_summary.csv:
% 29 of 125,401 (0.02%); results/eta_summary.csv: 41 of 67,608 (0.06%); results/limits_summary.csv: 0 of 571;
% results/large_k/large_k_summary.csv: 0 outside SG
Apart from SG with $100$ or more names, this happened in under $1\%$ of the decisions of every grid of real-data runs.

\begin{algorithm}[t]
\caption{One allocation decision after the close of day $t$.}
\label{alg:allocation}
\begin{algorithmic}[1]
  \REQUIRE Closes through day $t$; held shares $n^{\mathrm{prev}}$; wealth $W$ (USD); weighting rule $U$; selection
    rule; hedge on or off; whether a pick of the basket is due (Algorithm~\ref{alg:mpc}).
  \ENSURE Target shares $n^{\star}$ and orders for the open of day $t+1$, or the kept book and no orders.
  \STATE $\gamma_i \gets$ Eq.~\ref{eq:gamma} for every universe name, from the day's $(\alpha_i, \beta_i)$
    (Eq.~\ref{eq:ewls}), $\tilde g_M$ and $\xi_t$ (Eq.~\ref{eq:xi})
  \IF{a pick is due}
    \STATE $\mathcal{K} \gets$ picks of the selection rule; $\mathcal{H} \gets$ hedge pool if the hedge is on, else
      $\emptyset$
    \STATE $\mathcal{S}^{-} \gets$ up to $3$ names per sector with the most negative $\gamma_i < 0$, outside
      $\mathcal{H}$ \hfill \COMMENT{LS only; else $\emptyset$}
  \ENDIF
  \STATE $\mathcal{N} \gets \mathcal{K} \cup \mathcal{H}$; for LS,
    $\mathcal{N} \gets \{ i \in \mathcal{K} : \gamma_i > 0 \} \cup \mathcal{H} \cup
    \{ i \in \mathcal{S}^{-} : \gamma_i < 0 \}$
  \STATE Build $\mu_{g,i} = \alpha_i + \beta_i \tilde g_M$, $\Sigma_g$ (Eq.~\ref{eq:mpc-cov}), $\bar w_i$,
    $\bar c\,\bar S$, $B = 0.10\,W$, and $\bar\beta = \max(0, 1 - \xi_t)$ (hedge on) or $\bar\beta = \infty$
  \IF{$U$ uses $\gamma$ and no $i \in \mathcal{N}$ has $\gamma_i > 0$}
    \STATE $n^{\star} \gets$ floors: $\varepsilon$ (CD, CES, CESw), $0$ (LS), or $b_i$ plus the least hedge weight
      that meets $\bar\beta$ (SG) \hfill \COMMENT{no utility solve}
  \ELSE
    \STATE Solve Eq.~\ref{eq:allocation} for $U$ with Clarabel \hfill \COMMENT{MinVar, Tangent: mix, then scale}
    \IF{Clarabel is neither optimal nor almost optimal}
      \STATE Solve the same program with SCS
    \ENDIF
    \IF{a hedge name has $\gamma_i \le 0$ and $U$ uses $\gamma$}
      \STATE Fix the other names' shares; minimize $\sum_{i \in \mathcal{H}} w_i$ under the same limits
    \ENDIF
  \ENDIF
  \IF{the program is infeasible, both solvers failed, or some $n^{\star}_i$ is not finite}
    \STATE $n^{\star} \gets n^{\mathrm{prev}}$; \textbf{return} no orders \hfill \COMMENT{keep the whole book}
  \ENDIF
  \STATE $n^{\star}_i \gets 0$ where $\lvert n^{\star}_i S_i \rvert \le 0.001\,W$, and for held names outside
    $\mathcal{N}$
  \RETURN orders $n^{\star} - n^{\mathrm{prev}}$, filled at the open of day $t+1$
\end{algorithmic}
\end{algorithm}
